% \vspace{-4mm}
\section{Related work}
% \vspace{-3mm}
\noindent \textbf{Safety guidelines} \quad
Designing tasks that elicit unsafe behavior from AI agents requires grounding in established risk taxonomies and policies. Frameworks such as the AIR taxonomy~\citep{zeng2024airiskcategorizationdecoded} and technical interpretations of the EU AI Act~\citep{guldimann2025complaiframeworktechnicalinterpretation} define categories spanning operational, societal, and legal risks. Recent work emphasizes aligning agent behavior with human values~\citep{tang2024prioritizingsafeguardingautonomyrisks} and constructing environments that provide safe interaction affordances~\citep{chan2025infrastructureaiagents}. These perspectives inform the risk categories and scenario designs used in \OpenAgentSafety.

\noindent \textbf{LLM and agent safety evaluations} \quad
Prior benchmarks have focused extensively on unsafe generations from LLMs~\citep{röttger2025safetypromptssystematicreviewopen, tedeschi2024alertcomprehensivebenchmarkassessing}, probing biases, toxic completions, and jailbreaking strategies~\citep{doumbouya2025h4rm3llanguagecomposablejailbreak, jiang2024wildteamingscaleinthewildjailbreaks}. While these efforts helped shape safety-aligned finetuning and refusal training~\citep{kumar2023certifying, wang2023self}, they primarily assess static output generation. In contrast, agent safety works assess agents with tool-use capabilities~\citep{mo2024tremblinghousecardsmapping, li2025commercialllmagentsvulnerable}, expanding the risk surface to include execution-based harms. However, many such evaluations rely on simulated APIs and simplified environments~\citep{andriushchenko2025agentharmbenchmarkmeasuringharmfulness, yin2025safeagentbenchbenchmarksafetask, yuan2024rjudgebenchmarkingsafetyrisk}, limiting realism. Other evaluations are constrained to single tools and short interactions. Tool-specific evaluations have largely targeted:
(i) Web environments: Testing agents’ robustness to pop-ups, authentication barriers, and misleading content~\citep{tur2025safearenaevaluatingsafetyautonomous, zhang2024attackingvisionlanguagecomputeragents, xu2024advwebcontrollableblackboxattacks, chen2025shieldagentshieldingagentsverifiable};
(2) Code execution: Evaluating safety in generating or running scripts~\citep{guo2024redcoderiskycodeexecution}; and
(3) Social interaction: Simulating user conversations or agent collaboration~\citep{shao2025collaborativegymframeworkenabling, zhou2024sotopiainteractiveevaluationsocial}. Our work differs by integrating real tools (e.g., code execution, browsers, messaging) into a single framework with multi-turn, multi-user interactions. Unlike prior work, we simulate both benign and adversarial users, exposing agents to more realistic decision-making challenges.

\noindent \textbf{Training for safer agents} \quad
To improve agent robustness, recent work proposes scoring actions as safe or unsafe~\citep{yuan2024rjudgebenchmarkingsafetyrisk}, defensive agent architectures~\citep{chen2025shieldagentshieldingagentsverifiable}, and adversarial fine-tuning strategies~\citep{rosser2025agentbreedermitigatingaisafety}. Others advocate for active learning to prioritize rare risk cases~\citep{abdelnabi2025firewallssecuredynamicllm}, or explore how performance optimization can reduce safety margins~\citep{wu2025dissectingadversarialrobustnessmultimodal}. While promising, these approaches often assume access to evaluation settings that mirror realistic threats. Our benchmark fills this gap by offering a high-fidelity simulation framework suitable for safety training, adversarial red-teaming, and reinforcement learning setups.

